\section{\Dargent}{\label{sec:dargent}}

We have designed and implemented a data layout description language called \dargent,
which describes how a \cogent algebraic data type may be laid
out in (heap) memory, down to the bit level. Layout descriptions in \dargent are transparent to
the shallow embedding of \cogent's semantics, but they
influence the definition of the refinement relation to C code
generated by the compiler. In \autoref{sec:verification}, we describe in more
detail the extensions to our verification framework to accommodate
the \dargent layout descriptions. Here, we focus on the language
definition. Firstly, we give an informal overview of \dargent's language features.

\subsection{An Informal Introduction to \dargent} \label{ssec:informal-intro}

%%
%% CMCL: Reword; type vs. object distinction
%%
\Dargent offers the possibility to assign any boxed \cogent type a
custom layout describing how an object of that type should be stored in
the heap.
%
A boxed type
assigned with a custom layout is compiled to a C struct with a single
field: an array of 32-bit words.\footnote{Throughout the paper, unless
we explicitly specify the size, the term
``word'' always refers to unsigned integer types of 1 byte, 2 bytes,
4 bytes or 8 bytes and the actual size is usually less relevant in
the discussion. It does not necessarily imply pointer-sized words.
We discuss the implications
of the choice of the word size later in \autoref{ssec:gsetters}. 
}
This array \emph{represents} the
\cogent type: It can be deemed as untyped in C, but the \dargent description
contains enough information to access individual parts of the type correctly.
How data is laid out in memory is purely a low-level concern, and it does not
affect the functional semantics of a \cogent program in any way.
In other words, \cogent functions are parametric over the layouts of types.
%
Under the hood, the \cogent compiler generates custom getters (and
setters) in C, retrieving (and setting) the relevant parts of the object
from its array representation.

As an example, consider the \cogent type in \autoref{fig:cogent-type-ex}.
\begin{figure}
  \centering
  $$
  \begin{array}{ll}
    \mathbf{type}\ \AbsN{Example} = \{ \\
  \quad \FieldTy{struct}{\#\{\FieldTy{x}{\PrimType{U32}}, \FieldTy{y}{\PrimType{Bool}}\}},
    & \!\!\textit{-{}- nested embedded record} \\
  \quad \FieldTy{ptr}{\{\dots\}},
    & \!\!\textit{-{}- pointer to another record} \\
  \quad \FieldTy{sum}{\VariantTy{\Cons{X}{\PrimType{U16}} \alt \Cons{Y}{\PrimType{U8}}}}
    & \!\!\textit{-{}- variant type} \\
    \} 
  \end{array}
  $$
  \caption{\Cogent type example}
  \label{fig:cogent-type-ex}
\end{figure}
This record consists of three fields: $\FieldN{struct}$, $\FieldN{ptr}$ and $\FieldN{sum}$.
The type of the $\FieldN{struct}$ field is an unboxed record, denoted by the leading $\#$ symbol. This field is
embedded inside the parent $\AbsN{Example}$ record. The $\FieldN{ptr}$ field is a boxed
record, and is stored somewhere else in the heap, referenced by a pointer in the $\AbsN{Example}$
type. The last field $\FieldN{sum}$ has a variant type, with two alternatives tagged $\ConsN{A}$ and $\ConsN{B}$
respectively. This variant is unboxed (recall that there is no boxed variant in \cogent) and
is stored in the heap inside the parent record.

A layout for this record type must specify where each field is located in the
word array. Overlapping is not allowed, except for the payloads of the two constructors
of the variant type, since only one of them is
relevant at each time, depending on the tag value. The layout must also specify 
what the tag values for the variant constructors $\ConsN{A}$ and $\ConsN{B}$ are.
We will give a layout to this type after a short
introduction to the \dargent language constructs.

% Layouts cannot currently be assigned to \emph{unboxed} records, i.e., records that
% are not behind a pointer (for example, when declaring a record as a local
% variable). We plan to overcome this limitation in the future.
%

\begin{figure}
$$
  \begin{array}{llclr}
  \text{Sizes} & s & ::= & \multicolumn{2}{l}{n\rB \mid n\rb \mid s + s} \\
  \text{Layout expressions} 
     & \ell & ::=  & s & \text{(block of memory)}\\
     &      & \mid & x & \text{(layout variable)} \\
     &      & \mid & \mathsf{pointer} & \text{(pointer layout)}\\
     &      & \mid & L\ \many{\ell_i} & \text{(another layout)}\\
     &      & \mid & \ell\ \mathsf{at}\ s & \text{(offset operator)}\\
     &      & \mid & \ell\ \mathsf{after}\ \FieldN{f} & \text{(relative location)}\\
     &      & \mid & \ell\ \mathsf{using}\ \omega & \text{(endianness)}\\
     &      & \mid & \multicolumn{2}{l}{\mathsf{record}\ \{ \overline{\FieldN{f}_i : \ell_i} \}} \\
     &      & \mid & \multicolumn{2}{l}{\mathsf{variant}\ (\ell)\ \{ \overline{\ConsN{A}_i\ (n_i) : \ell_i} \}} \\
     % &      & \mid & \multicolumn{2}{l}{\mathsf{array}\ \{ \ell \}} \\
  \text{Declarations} & d & ::= & \multicolumn{2}{l}{\mathbf{layout}\ L\ \many{x_i} = \ell} \\
  \text{Layout names} &   & \ni & L\\
  \text{Endianness} & \omega & ::= & \bigEnd \mid \littleEnd \\
  \text{Natural numbers} &   & \ni & n, m\\
%  \text{Layout contexts} & C & ::= & \many{\ell_i \sim \tau_i} \\
\end{array}
$$
\caption{The syntax of the \Dargent surface language}
\label{fig:dargent-syntax}
\end{figure}


In \autoref{fig:dargent-syntax} we present the surface syntax for \Dargent.
A \emph{layout expression} is a description of the usage of some (heap) memory.
It only describes the low-level view of a memory region---it is not associated
to any particular algebraic data type. From this perspective, \dargent descriptions
are independent of \cogent types. As we will shortly see, however, a given \cogent type can only be laid out in certain ways, which places restrictions on which layouts can be assigned to a given type. 

When specifying a layout, two pieces of information are relevant: How much space a component
occupies in memory and where it is placed in relation to the overall heap object in which it is contained.
Primitive types, such as integer types and
booleans, are laid out as a contiguous block of memory of a particular
size. For example, a contiguous 4-byte block would be an
appropriate layout for the 32-bit word type $\PrimType{U32}$.
A block of memory is specified as a size expression,
which can be in bytes ($\byte$), bits ($\bit$), or additions of smaller sizes.
Additionally, memory blocks of word size (e.g.\ 1 byte, 2 bytes, 4 bytes and 8 bytes)
can be given an endianness ($\bigEnd$ or $\littleEnd$), with the $\using$ keyword.

Components of boxed type are represented as pointers, and thus must be described with the special $\DLPtr$ layout, and not as a chunk of
memory. This special layout improves readability of code, and also allows for some portability:
The $\DLPtr$ layout will have different sizes according to the host machine's architecture.


Layouts for record types use the $\mathsf{record}$ construct, which
contains sub-expressions for the memory layout of each field. As
we can specify memory blocks down to the individual bits, we can naturally
represent records of boolean values as a bit field:
%
\begin{mathpar}
  \mathbf{layout}\ \AbsN{Bitfield} = \DLRec{
     \FieldN{x}: 1\bit
   , \FieldN{y}: 1\bit\ \at\ 1\bit
   , \FieldN{z}: 1\bit\ \at\ 2\bit
  }
\end{mathpar}
%
Here the $\at$ operator is used to place each field at a
different bit offset, so that they do not overlap. If two record fields
have overlapping layouts, the description is rejected by the compiler.

The $\at$ operator can be applied to any layout expressions, which will
shift the entire expression by the specified amount. Alternatively,
the $\after$ operator can be used in a record layout:
\begin{mathpar}
  \mathbf{layout}\ \AbsN{Bitfield} = \DLRec{
     \FieldN{x}: 1\bit
   , \FieldN{y}: 1\bit\ \after\ \FieldN{x}
   , \FieldN{z}: 1\bit\ \after\ \FieldN{y}
  }
\end{mathpar}
so that a later field is placed right after the previous one. This saves
the programmer from calculating the concrete offset. When no offset ($\at$ or $\after$) is given,
it will by default place the field after the previous one.


Layouts for variant types use the $\mathsf{variant}$ construct. It
firstly requires a layout expression for the \emph{tag}. Then, for
each constructor in the variant, a specific tag value needs to be assigned, followed by
a layout expression for the \emph{payload} of that constructor. When one alternative is taken, the memory used by another alternatives becomes irrelevant, which is why the memory for the payloads can overlap.
Additionally, \dargent allows for zero-sized payloads. For instance,
the $\Cons{Maybe}{a}$ type, defined as $\VariantTy{\Cons{Just}{a} \mid \Cons{Nothing}{()}}$,
may be given a layout in which the payload for constructor $\ConsN{Nothing}$ does not occupy
any memory.


% \review{If you encode `Maybe T` as a tag followed by 7 bytes, can you then configure the compiler to replace addition on these `Maybe`s by an optimized computation that leverages the layout choice (e.g. `and`s the flags, performs native byte addition, and then resets the flags)?}
% 
% \ambroise{that could be a nice example of FFI. I'll implement it and describe it if I have time.}

% Array layouts are used for specifying the built-in array types in \Cogent.
% Currently we only allow a uniform layout for all elements of an array, thus
% in the $\mathsf{array}$ layout expression, we only need to supply a layout
% for the element type.

Similar to \cogent types, \dargent expressions are also structural. Layout synonyms
can be defined using the \texttt{layout} keyword, just as
\cogent type synonyms are defined using the \texttt{type} keyword. For example,
%
\begin{mathpar}
\mathbf{layout}\ \AbsN{FourBytes} = 4\byte
\end{mathpar}
%
defines a layout synonym $\AbsN{FourBytes}$, which is definitionally equal to $4\byte$ on
the right hand side. Layout and type synonyms can take parameters.

We can now give a \dargent description to the $\AbsN{Example}$ type in \autoref{fig:cogent-type-ex} (assuming a 64-bit architecture):
$$
\begin{array}{l}
  \mathbf{layout}\ \AbsN{ExampleLayout} = \mathsf{record}\ \{ \\
  \quad \FieldTy{struct}{\mathsf{record}\ \{\colorbox{red!10}{$\FieldTy{x}{4\byte}$}, \colorbox{yellow!30}{$\FieldTy{y}{1\bit}$}\}}, \\
  \quad \colorbox{cyan!10}{$\FieldTy{ptr}{\DLPtr\ \at\ 8\byte}$}, \\
  \quad \FieldTy{sum}{\mathsf{variant}\ (\colorbox{gray!30}{$1\bit$}) \\
  \quad \quad \{\colorbox{blue!10}{$\ConsN{X}(0) : 2\byte\ \at\ 1\byte$}, 
                \colorbox{green!10}{$\ConsN{Y}(1) : 1\byte\ \at\ 1\byte$}}\}\ \at\ 5\byte \\
  \} 
\end{array}
$$

\begin{figure}
\begin{tikzpicture}[background rectangle/.style={fill=white}, show background rectangle,scale=0.80]
  \draw [dashed] (0,0) rectangle (16,2) ;
  \draw [dotted,fill=gray!10] (5,0.1) rectangle (8,1.9) ;
  \draw [fill=red!10] (0.1,0.1) rectangle (4,1.9)node [pos=0.5,align=center] {$\FieldN{x}$};
  \draw [fill=yellow!30] (4.00,0.1) rectangle (4.125,1.9)node [pos=0.5,align=center,anchor=west] {$\FieldN{y}$};
  \draw [fill=gray!30] (5,0.2) rectangle (5.125,1.8);
  \draw [fill=blue!10] (6,1.2) rectangle (8,1.8)node [pos=0.5,align=center] {$\ConsN{X}$};
  \draw [fill=green!10] (6,0.2) rectangle (7,0.8)node [pos=0.5,align=center] {$\ConsN{Y}$};
  \draw [fill=cyan!10] (8,0.1) rectangle (15.9,1.9)node [pos=0.5,align=center] {$\FieldN{ptr}$};
  \draw [->] (5.125,1.5) -- (6,1.5) node [pos=0.5, anchor=south, yshift=-2pt] {\scriptsize if $\!= \!0$};
  \draw [->] (5.125,0.5) -- (6,0.5) node [pos=0.5, anchor=south,yshift=-2pt] {\scriptsize if $= 1$};
  \node [align=center, anchor=north] at (0,0) {0B};
  \node [align=center, anchor=north] at (1,0) {1B};
  \node [align=center, anchor=north] at (2,0) {2B};
  \node [align=center, anchor=north] at (3,0) {3B};
  \node [align=center, anchor=north] at (4,0) {4B};
  \node [align=center, anchor=north] at (5,0) {5B};
  \node [align=center, anchor=north] at (6,0) {6B};
  \node [align=center, anchor=north] at (7,0) {7B};
  \node [align=center, anchor=north] at (8,0) {8B};
  \node [align=center, anchor=north] at (9,0) {9B};
  \node [align=center, anchor=north] at (10,0) {10B};
  \node [align=center, anchor=north] at (11,0) {11B};
  \node [align=center, anchor=north] at (12,0) {12B};
  \node [align=center, anchor=north] at (13,0) {13B};
  \node [align=center, anchor=north] at (14,0) {14B};
  \node [align=center, anchor=north] at (15,0) {15B};
  \node [align=center, anchor=north] at (16,0) {16B};
\end{tikzpicture}
  \caption{The $\mathit{ExampleLayout}$, visualised.}
  \label{fig:example-layout}
\end{figure}
\noindent \autoref{fig:example-layout} gives a pictorial illustration of this memory layout.
In the layout above, it is worth noting that the offsets ($\at\ 1\byte$) for the two payloads of $\ConsN{A}$ and $\ConsN{B}$
are in relation to the beginning of the $\FieldN{sum}$ field, which is $5$ bytes ($5\byte$) from the beginning
of the top-level structure. We can equivalently write $\DLVar{1\bit\ \at\ 5\byte}{\ConsN{A}(0) : 2\byte\ \at\ 6\byte,
\ConsN{B}(1) : 1\byte\ \at\ 6\byte}\ \at\ 0\byte$ for the $\FieldN{sum}$ field without changing its layout.

At this point, this layout is still independent of the \cogent type, and the compiler will only check
that this layout definition is wellformed: that it does not have overlapping fields, that the tag values
are distinct, and so on. To associate a layout to a type, we add a $\layout$ keyword to the type language of \cogent. For this example, the type $\AbsN{Example}$ $\layout$ $\AbsN{ExampleLayout}$ describes the type $\AbsN{Example}$ laid out according to the description in $\AbsN{ExampleLayout}$. The compiler will check that $\AbsN{ExampleLayout}$ is an appropriate layout for the \cogent type $\AbsN{Example}$. We will talk more about the wellformedness and matching rules later
in \autoref{ss:formalised-semantics}. To eliminate verbosity, a type synonym can be given to the layout-annotated type above.


We make the design choice of having \dargent layouts be defined independently 
of \cogent types, and only relating them afterwards with the $\layout$
keyword. This design may seem sub-optimal, as some common information is duplicated
in both the types and the layouts, and matching them requires a set of
dedicated typing rules (as we will
see in \autoref{fig:layout-type}). We make this design decision for several reasons.

Firstly, while the \cogent type structure is central to the functional semantics of any \cogent
program and reasoning about its functional correctness, the exact layout of algebraic data types is just an
implementation detail, totally irrelevant to the top-level Isabelle/HOL
embedding of the \cogent program, and whose correctness is guaranteed by the automatic
C refinement proof that the \cogent compiler produces. 
%
Thus, type structure and layout are conceptually separate.
%These two constructs are conceptually separate.

Secondly, this approach is more flexible and amenable to future
extensions. Currently, as we will see in \autoref{fig:layout-type}, the
layout--type matching is fairly restricted: e.g.\ a $\PrimType{U16}$
type has to occupy 2 bytes ($2\byte$) in memory, and a record type has
to be laid out in accordance with a $\mathsf{record}$ layout.  In the
future, several extensions could be made to relax this matching
relation. For example, it is technically valid to store a smaller type
in a larger memory area, say, a $\PrimType{U16}$ type in a memory region
of 4 bytes, or a record type in a contiguous chunk of memory that is big
enough. Some heuristics can be implemented in the compiler to decide how
to arrange underspecified layouts. This feature can be useful in
improving programmers' productivity, permitting the omission of layout
details for unimportant parts of a type. Also, as there is ongoing work
towards adding refinement types to \cogent~\citep{Paradeza_20}, a
refined type could possibly be laid out in a smaller memory area. For
instance, $\{ \nu : \PrimType{U16}\ |\ \nu < 2 \}$ can be laid out in a
memory area as small as 1 bit. None of these extensions would be easy to
implement if the layout information was baked into the types.

Thirdly, separating layouts and types encourages modularity.
Developers using the \cogent language can write programs without
needing to worry about the detailed layout of types,
and are still able to prove functional correctness of their code
against some high-level specification and ship their code to end users.
The end users, with the knowledge of the particular target architecture and environment,
can decide the layouts and plug them into the \cogent programs.

Finally, although our design requires a dedicated set of rules for
checking the layout--type matching, it significantly simplifies compiler
engineering in the long term. The \cogent compiler is very large, and
the layout--type matching checker only constitutes a small part of the
typechecker. The compiler not only compiles \cogent programs to C, but
also generates information for various Isabelle proof tactics, numerous
embeddings of the program in Isabelle and Haskell~\citep{Chen:22}. A lot
of these embeddings are only concerned with types, and not layouts. If
the layouts and types were merged, any changes to the layout
implementation would require changes to various irrelevant parts of the
compiler.


\subsection{Layout Polymorphism}

% \dargent layouts are only concerned of low-level implementation details of a \cogent program,
% and they do not affect the functional semantics of a program. It means that program semantics
% is parametric over layouts.
We extend \Cogent's existing parametric polymorphism mechanism to support \emph{layout polymorphism}. This feature allows users to abstract
over the layout that they would like to assign to a certain type.
In a \cogent function signature, \emph{layout variables}, just as type variables, can
be universally quantified. These layout variables may be constrained, similarly to type constraints in Haskell, which require that a layout variable matches a type.\footnote{A formal definition of layout matching is discussed in \autoref{ss:formalised-semantics}.}
For example, in the code snippet below,
\begin{align*}
  & \mathbf{type}\ \Cons{Pair}{t} = \RecordTy{\FieldTy{fst}{\VarN{t}}, \FieldTy{snd}{\VarN{t}}} \\
  & \layout\ \Cons{LPair}{l} = \DLRec{\FieldTy{fst}{\VarN{l}}, \FieldTy{snd}{\VarN{l}\ \at\ 4\byte}}\\
  & \FunN{freePair} : \forall (t, l :\sim t).\; \Cons{Pair}{t}\ \withL{\Cons{LPair}{l}} \rightarrow ()
\end{align*}
we define a parametric type synonym $\Cons{Pair}{t}$ and layout synonym $\Cons{LPair}{l}$, and
an abstract polymorphic function that operates on such a pair. In the function's type, we require layout $l$ to be compatible with type $t$, so that the $\Cons{LPair}{l}$
is always a valid layout expression for type $\Cons{Pair}{t}$.
Layout-polymorphic functions may be explicitly applied to layouts, akin to explicit type applications.
For example, $\TyLyApp{freePair}{\PrimType{U32}}{\AbsN{FourBytes}}$ instantiates $t$ to $\PrimType{U32}$
and $l$ to $\AbsN{FourBytes}$. If the type/layout application is
incompatible, as for instance in $\TyLyApp{freePair}{\PrimType{U8}}{1\bit}$, where
$1\bit$ is not big enough to store the $\PrimType{U8}$, the typechecker will reject
such a program.  The typechecker also ensures that any instantiation of $l$ produces layouts
that are wellformed. For example, if $l$ is instantiated with $8\byte$, it will
render $\Cons{LPair}{l}$ ill-formed, as the $\FieldN{snd}$ field will overlap with the
$\FieldN{fst}$. This can be rectified by using the $\after$ relative location (or
leaving the location implicit)
instead, which will automatically place the $\FieldN{snd}$ field right after $\FieldN{fst}$.

Layout polymorphism is necessary to retain the full generality of type
polymorphism in the presence of \dargent, as demonstrated by the example
above.
%
Layout polymorphism also facilitates code reuse in several scenarios:
\begin{enumerate*}
  \item It can be used in programs that are architecture-agnostic.
    For the same program running on different architectures, the same algebraic datatype
    may be laid out differently according to the hardware it runs on. These datatypes are
    not necessarily part of the API, as developers tend to design interfaces
    in an architecture-independent manner and these types thus have uniform layouts
    across different architectures. However, datatypes that are internal to the program
    can be represented with appropriate layouts to best suit the architecture and the
    respective C compiler for better performance.
  \item Besides architectural differences, layout polymorphism allows types to be reused more
    broadly across different applications.
    For instance, colours are a general concept independent of the application. One common
    way to logically represent a colour is the ARGB model. However, the actual low-level layouts
    for the logical ARGB value vary. Commonly used layouts include ARGB32 and RGBA32, and
    other layouts also exist. In this case, developers can simply define an ARGB colour type with a polymorphic layout to be determined at use-sites.
  \item Layout polymorphism also facilitates software engineering practices.
    For example, developers can start with the default layout without a \dargent annotation,
    and \emph{incrementally}
    optimise the program to use more compact and efficient \dargent layouts. In this
    case, they simply instantiate the layouts differently, without needing to
    re-define the types. Benchmarking is another use case: developers can easily have
    the same type with different instantiations \emph{side-by-side} to study their performance
    differences.
\end{enumerate*}


%In Section~\ref{ssec:bitfields}, we present another extension that accounts for custom sized integers such as $\ConsN{U7}$ or $\ConsN{U15}$. \liam{Don't understand why this is mentioned here}

%Additionally, choosing to keep layouts and types separate does not preclude us from providing users with syntactic sugar that eliminates information duplication. 
% implementing syntactic sugar
%to eliminate information duplication
% and improve the user
%experience of the language. 
%We leave it
%as future work.

\subsection{The \Dargent Core Calculus}\label{ss:dargent-core}

\begin{figure}
\[
\begin{array}{llclr}
  \text{Bit ranges} & r\\
  \text{Offsets}    & o & \in & \mathbb{N} \\
  \text{Layouts} & \ell & ::=  & () & \text{(unit layout)}\\ 
                 &      & \mid & r_\omega & \text{(primitive layout)} \\
                 &      & \mid & x\{o\} & \text{(layout variable)} \\
                 &      & \mid & \multicolumn{2}{l}{\mathsf{record}\ \{ \overline{\mathrm{f}_i\!: \ell_i} \}}\\
                 &      & \mid & \multicolumn{2}{l}{\mathsf{variant}\ (\ell)\ \{ \overline{\mathrm{A}_i\ (n_i)\!: \ell_i} \}}\\
 %                &      & \mid & \multicolumn{2}{l}{\mathsf{array}\ \{ \ell \}}\\
  \text{Field names} & & \ni & \mathrm{f}\\
  \text{Constructors} & & \ni & \mathrm{A}\\
  \text{Endianness} & \omega & ::= & \multicolumn{2}{l}{\mathsf{BE} \mid \mathsf{LE} \mid \mathsf{ME}} \\
%  \text{Layout contexts} & C & ::= & \many{\ell_i \sim \tau_i} \\
\end{array}
\]
\caption{The syntax of the \Dargent core language}
\label{fig:dargent-core}
\end{figure}




The \Dargent surface language is desugared into a smaller core calculus,
whose syntax is outlined in \autoref{fig:dargent-core}. The core calculus is
the language on which the verification is based. As can be seen in  \autoref{fig:dargent-core}, layouts consist fundamentally of \emph{bit ranges}, which describe which bits in memory are used
to store each piece of data. The definition of bit ranges is omitted, because our core calculus is parametric over the exact bit range
representation---at different points in the compilation pipeline, bit ranges are represented
differently, to suit the purpose of that particular phase of the compilation.

Bit ranges are annotated with an endianness to form a primitive layout in our 
core language. When the endianness is not specified in the surface language, 
a \emph{machine endianness} will be given by default, written $\mathsf{ME}$
 in core. When C code is generated, a
subroutine will be invoked to determine the machine endianness, so that the C code works
as intended.
%
In this paper, when the endianness is unimportant, we will omit the
subscript from $r_{\omega}$.

When the surface layout expressions are first desugared, the bit ranges are represented
as $\bitrange{o}{s}$, which is a pair of natural numbers, indicating the offset ($o$ bits) from
\emph{the beginning of the top-level heap object} in which it is contained, 
as well as the range occupied ($s$ bits).
%
In $\bitrange{o}{s}$, we require $s > 0$; zero-sized layouts can use the
empty layout $()$ instead.

Layout variables are given the form $x\{o\}$, since we must remember to
apply the offset after instantiation---when $\VarN{x}$
is instantiated, it will be shifted to the right by $\VarN{o}$ bits.
The other core layouts are very similar to their counterparts in the
surface language.

In summary, the desugarer converting the surface layout expressions to
the core calculus must perform the following tasks:
\begin{itemize}
  \item expand layout names to layout definitions;
  \item translate size expressions into bit ranges;
  \item compute offsets relative to the beginning of the top-level
    heap object;
  \item insert explicit layout applications to any layout-polymorphic
    function calls.
\end{itemize}
% The desugaring algorithm is a simple inductive definition, which can be 
% found in~\autoref{ss:static-semantics}.
\Dargent is used for describing the layout of heap memory. To cover all the
heap memory that is addressable from an object,
it suffices to attach a layout description to each pointer that the object contains.
In the $\AbsN{Example}$ type that we showed earlier in \autoref{ssec:informal-intro},
the layout annotation $\AbsN{ExampleLayout}$ is attached to the $\AbsN{Example}$ type,
which is a boxed type. That layout description contains all the information about
how the fields should be laid out, in particular the unboxed fields $\FieldN{struct}$ and
$\FieldN{sum}$ which also reside in the heap. For the boxed field $\FieldN{ptr}$, however,
the description in $\AbsN{ExampleLayout}$ does not extend beyond the indirection---it only
knows that $\FieldN{ptr}$ is a $\DLPtr$ but it is oblivious to the
memory layout behind $\FieldN{ptr}$.
%
To prescribe the layout for $\{ \FieldN{c} : \PrimType{U8}\}$, a
separate layout annotation should be attached to this type, \emph{e.g.}
$\{ \FieldTy{c}{\PrimType{U8}} \}\ \layout\ \DLRec{
  \FieldTy{c}{1\byte}}$.

Recall that, if an object is referenced by-pointer, it has a boxed type.
This means that in the core calculus, data layouts only need to be attached to boxed sigils
and nothing else.
We modify the definition of sigils as follows:
$$
\begin{array}{llclr}
  \text{Sigils} & s & ::=  & \Writable{\ell} \mid \ReadOnly{\ell} \mid \Unboxed &
\end{array}
$$
If the sigil is boxed (either writable or read-only), it can carry a \dargent layout.
We use $\Boxed{\ell}$ as a notational convenience when we do not wish
to distinguish its mutability.
When the layout annotations are left out, the type will be compiled with the default layout chosen
by the compiler.



\subsection{The Static Semantics with \Dargent} \label{ss:formalised-semantics}

% \review{any proof of the soundness of the type system?}

%type system, kinding system for Cogent
%changes needed for Dargent both to type and kinding system
% we need to ensure that types and layouts match.
% Since cogent has polymorphic types and layouts. 
% We do not always have sufficient information at the type checking phase to check that types and layouts match. 
% we therefore need to track a set of pairs of types and layouts that we need to check matching for at a later stage of compilation (after monomorphisation).

% lemma 3 in cogent paper 
%Lemma 3 (Type instantiation preserves kinds). For any type τ, αi:Kκi ⊢τ:K κimplies∆⊢τ[ρi/αi]:K κwhen,foreachi, ∆⊢ρi :K κi.
% would no longer hold once we introduce layouts
%why? because now typing and kinding also need to substitute layouts and 
% since typing should match layouts and types, we need to track type variables and their layouts to ensure once they are instantiated, their instatiation matches the given layouts. Layouts can also be polymorphic and the type-layout matching should also be respected when layouts are instantiated. 
%
%Since layouts are also polymorphic, we also need to instantiate them
%1) if substitution is wf and type is wf then resulting type is wf. (why? subject reduction?)
%% need to introduce a set of constraints C

% 2) if a type is a subtype (in a linear sense) of a wf type then it is also wf and vice-versa.
% now a type wfness includes constraints C
% any type that includes a polym. type variable or poly layout needs to be checked for wf later on once these are instantiated. We therefore need to keep track of such types and layouts in order tpo check later on that they match and respect layout wf.  
% suppose you have a record with a layout, the context must contain record type with layout for the layout to be wf. 
% record type t with layout l (C must contain pair (t, l) )
% now we want to know that (bang(t), l) is also wellformed (why?)
% so constraints C must contain (bang (t), l)
%  linearity of types does not affect their layout in memory.
% i.e. whether types are read-only or writable
%therefore the matching between types and layouts can use simplified types (with linearity info removed)
% in fact this is necessary, why?
%  any type (regardless of linearity ) relates to layout
% because we want 3
% 3) if a type t is wf then the bang (t) is wellformed (turning to non-linear read only version)
% \delta = [\tau_1.. tau_n]
% [\tau:= \delta \tau], 

%Thereofore, Dargent's typing judgement needs to therefore additionally track 

% can't do that directly but rather need to keep track of a set of pairs (l,t) where t is a simplified type.
%
%a set of (l,t), the typesystem needs to ensure they match
% why simplified ?


%This section, we give a formal account of the static semantics of the core
%calculus on which we base our \dargent formalisation.

\Cogent's static semantics, along with some meta-properties
of the type system such as type preservation, are formalised in Isabelle/HOL
(\citet{OConnor:2016,Cogent_JFP_2021} give more details). In this section, we explain
how these formal proofs and definitions are extended to account for \Dargent.

In the original \Cogent typing judgement $\Delta;\Gamma \vdash e : \tau$, we have the set of polymorphic type variables $\Delta$ in the context. Each entry in $\Delta$ also includes a set of \emph{constraints} for that type variable, which may enforce, for example, a type variable represent only types that are shareable.
 When these type variables are instantiated, these constraints are checked. We also have a wellformedness judgement on types, $\Delta \vdash \tau\ \mathbf{wf}$, which ensures that all type variables are in scope in $\Delta$ and that the type structure of $\tau$ respects the linearity constraints in $\Delta$. 
%Given a \cogent type, \TODO{say something about subtyping}.  
%\Cogent's kinding judgement, $\Delta \vdash \tau :_K \kappa$, ensures that \cogent types are wellformed with respect to linearity constraints specified in the sigil.  
%$\Delta \vdash \tau :_K \kappa$, stating that the type $\tau$ is wellformed
%in the context of type variables $\Delta$. 

With \Dargent, the typing judgement must additionally ensure that \cogent types
match their assigned layouts (\autoref{fig:layout-type}) and that the layouts
are wellformed (\autoref{fig:wf-layout}). Because \cogent has
polymorphic type and layout variables, however, these additional constraints
cannot always be immediately ensured when typechecking locally. Thus, we
extend the \cogent typing judgement $\Delta;C;\Gamma \vdash e : \tau$ and
wellformedness judgement $\Delta; C \vdash \tau\ \mathbf{wf}$ to include an
additional set of \emph{layout} constraints $C$. It
is a set of pairs $(\ell,\hat{\tau})$ denoting that the layout $\ell$ must
be capable of representing type $\hat{\tau}$. When the types and layouts are instantiated, these matching constraints are checked.


\begin{figure}
\[
\begin{array}{llclr}
  \text{Simplified types} 
     & \hat{\tau} & ::=  & T \mid \mathsf{pointer} \mid () \mid \alpha \mid
                         \{ \many{\mathrm{f}_i : \hat{\tau}_i} \} \mid
                         \VariantTy{\many{\mathrm{A}_i\ \hat{\tau}_i}} \mid \dots &
\end{array}
\]
\caption{The syntax of simplified types}
\label{fig:syntax-simplified-types}
\end{figure}


The type in the context $C$ is \emph{simplified}, denoted as $\hat{\tau}$,
because any type structure beyond the single heap object described by the
layout $\ell$ is irrelevant. This means that all boxed types, which are all
uniformly represented as a pointer to another heap object, can be treated
simply as a primitive machine word when matching with layouts. The syntax of (a
subset of) simplified types is given in \autoref{fig:syntax-simplified-types}
(cf.\ \autoref{fig:syntax-cogent-types}). Types are simplified by
replacing all boxed components of the type with the newly introduced
$\mathsf{pointer}$ type constant.  

%
%
%
%Two important things to note at this stage: 
%(1) \cogent types can contain polymorphic type variables and \dargent supports polymorphic layouts. 
%(2) the layout of a type is independent of its linearity constraints
%
%In particular, the typing judgement in the presence of \dargent $\Delta;C;\Gamma\vdash e : \tau$ additionally track a set of pairs $C$ of layouts $\ell$ and \emph{simplified types} $\hat{t}$ (\autoref{fig:syntax-simplified-types}). \dargent's type system  maintains matching constraints between the types and the layouts. A simplified type is just a \Cogent type with the linearity information removed. 

%In this section, we give a formal account of the static semantics of the core
%calculus on which we base our \dargent formalisation.
%The typing rules of \dargent differ from those of the \Cogent  in that they involve an additional environment that lists constraints on relating a layout to a \emph{simplified type}
%(\autoref{fig:syntax-simplified-types}).


%A simplified type %
%\footnote{The notion of simplified types introduced here is very similar
%  to the \emph{representation types} used in the C refinement proof. The main difference is that
%  it now includes type variables, whereas the representation types don't, which makes sense as
%  C is not polymorphic. The properties that can be proved on this simplified types are also very similar
%  to those of representation types, especially stabilty under subtyping, which is a main
%  motivation and contribution of the subtyping design.}
%is just a \cogent type with all the linearity information removed. As such, the simplified type is stable under
%subtyping and $\mathbf{bang}(\cdot)$. 
%(turning a linear type into a read-only
%non-linear one), thus making constraints invariant
%under these operations. 
%These properties are used when adapting the proofs that subtyping and $\mathbf{bang}(\cdot)$ preserve wellformedness of types. 
%More precisely, the function
%$\textbf{simplify}$ that turns a regular \cogent type into a simplified type (boxed
%types are turned into pointers). Then,
%$\textbf{simplify}(\BangF{\tau}) = \textbf{simplify}(\tau) $ and
%$\textbf{simplify}(\tau_1) = \textbf{simplify}(\tau_2)$ if $\tau_1$ is a subtype
%of $\tau_2$.

\subsubsection{Typing Rule for Records with Layouts}

\newcommand{\highlight}[1]{\colorbox{blue!10}{$\displaystyle #1$}}
\newcommand{\highlighttext}[1]{\colorbox{blue!10}{#1}}

We are now ready to present the main typing rule that is added for
\dargent: wellformedness of boxed records with custom layouts.
\begin{equation*}
  \label{eq:wellformed-rec}
   \inferrule
	{\overline{\mathrm{f}_i}\ \text{distinct} \\
   \text{for each $i$: }\Delta; C \vdash \tau_i\ \mathbf{wf} \\\\
     \highlight{\ell\ \mathbf{wf}} \\
     \highlight{
       % \text{if $\ell$ or some $\tau_i$ is open, then}
       \ell\sim \{ \overline{\mathrm{f}_i :
         \hat\tau_i}\} }  
     \\
	 \highlight{
       % \text{if $\ell$ or some $\tau_i$ is open, then}
       (\ell,\{ \overline{\mathrm{f}_i :
         \hat\tau_i}\}) \in C}
   }
	{\Delta ; C\vdash \{ \overline{\mathrm{f}_i :: \tau_i} \} b_\ell\ \mathbf{wf}}
   {}
  {\textsc{RecWf}}
\end{equation*}
The first two premises in this rule are similar to those in \cogent, the \highlighttext{next three premises} are added to account for \dargent. The layout wellformedness judgement, $\ell\ \mathbf{wf}$, ensures that field layouts do not overlap (\autoref{fig:wf-layout}). 
The last two premises require the type to match the layout (\autoref{fig:layout-type}) and ensure that this is tracked in the set of constraints, respectively.
% \footnote{
%   The notion of context does not enforce that if $(\ell, \hat{t}) \in C$,
%   then $\ell \sim \hat{t}$, but 
%   \textsc{RecWf} is the only rule that requires some
%   constraint to be in $C$, and it indeed enforces it.
% }% 
In practice, we use a variant of the rule \textsc{RecWf} where the
layout constraint is only tracked in $C$ if the type or the layout are
polymorphic.



%These two premises are necessary for ensuring that any valid instantiation of a well-typed polymorphic program is well-typed which in turn is needed for proving type preservation. 
%the the following specialisation lemmas hold.
\begin{figure}
\begin{inductive}{\ell \sim \hat{\tau}}
  \inferrule
    {\mathrm{bits}\;(T) = s}
    {\bitrange{o}{s} \sim T}
    {\textsc{PrimTyMatch}}
  \qquad
  \begin{array}{ll}
    \bits{\ConsN{U}n}   & = n  \\
    \bits{\ConsN{Bool}} & = 1
  \end{array}
  \\
  \inferrule
    {M\ \text{is the pointer size}}
    {\bitrange{o}{M} \sim \mathsf{pointer}}
    {\textsc{BoxedTyMatch}}
  \qquad
  \inferrule
    { }
    {x\{o\} \sim \hat{\tau}}
    {\textsc{VarMatch}}
  \\
  \inferrule
    { }
    {() \sim ()}
    {\textsc{UnitMatch}}
  \qquad
  \inferrule
    {\text{for each $i$: } \ell_i \sim \hat{\tau}_i}
    {\DLRec{\many{\FieldN{f}_i : \ell_i}} \sim \{\many{\FieldN{f}_i : \hat{\tau}_i}\}}
    {\textsc{URecordMatch}}
  % \\
  % \inferrule
  %   {\ell \sim \tau \\
  %    \ell\ \text{is byte aligned}}
  %   {\DLArr{\ell} \sim \ArrayTy{\tau}{v}{\Unboxed}}
  %   {\textsc{UArrayMatch}}
  \\
  \inferrule
    {\text{for each $i$:\ } \ell_i \sim \hat{\tau}_i}
    {\DLVar{s}{\many{\ConsN{A}_i\ (v_i) : \ell_i}} \sim \VariantTy{\many{\ConsN{A}_i\ \hat{\tau}_i}}}
    {\textsc{VariantMatch}}
\end{inductive}
\caption{Matching relation between layouts and types}
\label{fig:layout-type}
\end{figure}



\begin{figure}
\begin{inductive}{\ell\ \mathbf{wf}}
  \inferrule
    { }
    {\bitrange{o}{s}\ \mathbf{wf}}
    {\textsc{BitRangeWF}}
    \qquad
    \inferrule
    { }
    {()\ \mathbf{wf}}
    {\textsc{UnitWF}}
    \qquad
    \inferrule
    { }
    {x\{o\}\ \mathbf{wf}}
    {\textsc{VarWF}}
  \\
  \inferrule
  {\text{for each $i$: } \ell_i \ \mathbf{wf} \\\\
    \text{for each $i\neq j$: }
    \mathrm{taken}(\ell_i) \cap 
    \mathrm{taken}(\ell_j) = \emptyset
  }
    {\DLRec{\many{\FieldN{f}_i : \ell_i}} \  \mathbf{wf}}
    {\textsc{URecordWF}}
  % \\
  % \inferrule
  %   {\ell \sim \tau \\
  %    \ell\ \text{is byte aligned}}
  %   {\DLArr{\ell} \sim \ArrayTy{\tau}{v}{\Unboxed}}
  %   {\textsc{UArrayWF}}
  \\
  \inferrule
  {\text{for each $i$:\ } \ell_i \ \mathbf{wf}\\
    v_i < 2^s
    \\
    \mathrm{taken}(\ell_i) \cap \mathrm{taken}(\bitrange{o}{s}) = \emptyset
    \\\\
    \text{for each $i\neq j$:\ } v_i \neq v_j }
    {\DLVar{\bitrange{o}{s}}{\many{\ConsN{A}_i\ (v_i) : \ell_i}}\ \mathbf{wf}}
    {\textsc{VariantWF}}
\end{inductive}

  where
  $\mathrm{taken}(\ell) \in \mathbb{N}$
   returns the set of bit positions that $\ell$ occupies
    (defined recursively).
\caption{Wellformed layouts}
\label{fig:wf-layout}
\end{figure}




\subsubsection{Specialisation}

% change substitution syntax to something more readable, remove kinding and just keep wellformedness of types if possible
\cogent's type preservation proof relies on type specialisation maintaining welltypedness~\citep[Lemmas 3 and 4]{OConnor:2016}. 
With \Dargent, we need to adapt these statements as follows.
%We state them below and discuss them afterwards, forgetting kinding information in $\Delta,\Delta'$ for the sake of simplicity.
\begin{lemma}[Specialisation lemmas]
  \label{lem:specialisation-lemmas}
  Let $\overline{\rho_i}$ be a list of types and $\overline{\ell_j}$ be a list
  of layouts. Let $\overline{\alpha_i}$ denote the list of type variables declared
  in $\Delta$.

  Then, $\Delta;C \vdash \tau\ \mathbf{wf} $ and  $\Delta;C;\Gamma \vdash e : \tau$  
  respectively imply 
     \begin{align*}
      \Delta';C'& \vdash \tau[\overline{\rho_i}/\overline{\alpha_i}, \overline{\ell_j}/\overline{x_j}]\ \mathbf{wf}; \\
      \Delta';C';\Gamma[\overline{\rho_i}/\overline{\alpha_i}, \overline{\ell_j}/\overline{x_j}]%^\dagger 
      & \vdash e[\overline{\rho_i}/\overline{\alpha_i}, \overline{\ell_j}/\overline{x_j}]: \tau[\overline{\rho_i}/\overline{\alpha_i}, \overline{\ell_j}/\overline{x_j}  ],
     \end{align*}
  under the following conditions:
  \begin{itemize}
    \item for each $i$, $\Delta', C'\vdash \rho_i\ \textbf{wf}$;
    \item for each pair $(\ell, \hat{\tau}) \in C$ such that
      $\Delta \vdash \hat{\tau}\ \mathbf{wf}$, the following statements hold:\footnote{The
        condition $\Delta \vdash \hat{\tau}\ \mathbf{wf}$ is required to support a stronger induction hypothesis. As always, the gory details are in the machine-checked proofs. }
      \begin{itemize}
      \item $\ell[\overline{\ell_j}/\overline{x_j}]\ \mathbf{wf}$;
      \item 
        $ \ell[\overline{\ell_j}/\overline{x_j}] \sim \hat{\tau}[\overline{\hat{\rho}_i}/\overline{\alpha_i}]$;
       \item $(\ell[\overline{\ell_j}/\overline{x_j}], \hat{\tau}[\overline{\hat{\rho}_i}/\overline{\alpha_i}]) \in C' $.
      \end{itemize}
  \end{itemize}
\end{lemma}

% \begin{lemma}[Type wellformedness specialisation]
%   \label{lem:type-specialisation}
%   If $\Delta;C \vdash \tau\ \mathbf{wf} $ and
%   $\Delta';C' \vdash \overline{\rho_i},\overline{l_j} :_s \Delta; C$, then
%   $$\Delta';C' \vdash \tau[\overline{\rho_i}/\overline{\alpha_i}, \overline{l_j}/\overline{x_j}] \mathbf{wf} $$
% \end{lemma}

\noindent In this lemma,  $[\overline{\rho_i}/\overline{\alpha_i}, \overline{\ell_j}/\overline{x_j}]$ denotes the simultaneous
substitution replacing both type variables $\alpha_i$ and layout variables $x_j$,
and $\Gamma$ is a typing context environment.


Informally, the two conditions mean that the pair 
$(\overline{\rho_i}, \overline{\ell_j})$ of type and layout lists defines a valid 
substitution from $\Delta;C$ to $\Delta';C'$. %, taking $\Delta := \overline{\alpha_i}$,
They also appear in the typing rule for specialising polymorphic functions.
The second one is new and enforces that for each  pair
$(\ell, \hat{\tau})\in C$ that is wellformed (in the sense that type variables
appearing in $\hat{\tau}$ are in $\Delta$), 
the substituted constraint is satisfied and remembered in the set of constraints $C'$.
This last requirement can in fact be dropped if the substituted constraint is closed.
%   Formally, the substitution wellformedness rule is defined as:

% $$
% \inferrule
%   { \text{for each}\ \kappa_i\ \text{in}\ \Delta\ \text{and}\ t_i\ \text{in}\ \delta:\ \Delta'; C' \vdash t_i :_K \kappa_i \\\\
%     \text{for each}\ (\ell_j, \hat{t}_j)\ \text{in}\ C:\ \ell_j[\varsigma] \sim \hat{t_j}[\textbf{simplify} \circ \delta] \\\\
%     (\ell_j[\varsigma], \hat{t_j}[\textbf{simplify}\circ \delta])\in C'
%   }
%   {\Delta'; C' \vdash \varsigma, \delta :_s \Delta; C }
%   {\textsc{SubstWf}}
% $$


% A similar specialisation lemma holds for well-typed terms as well:
% \begin{lemma}[Typing specialisation]
%   \label{lem:typing-specialisation}
%   If $\Delta;C;\Gamma \vdash e : t$ and
%   $\Delta';C' \vdash \varsigma,\delta :_s \Delta; C$, then
% 	$$\Delta';C';\Gamma[\delta, \varsigma]^\dagger \vdash e[\delta, \varsigma]: t[\delta, \varsigma].
% 	~\footnote{$\cdot[\cdot,\cdot]^\dagger$ is the simultanous substitution lifted to the environment.}
%   $$
% \end{lemma}

% Now, let us explain why the wellformedness rule \textsc{RecWf} for records
% not only requires a check that the layout matches the type, but also that
% this constraint is remembered. Indeed, because the type and the layout may
% contain variables, this matching is only partial and would relate
% the layout $ \mathsf{record}\ \{ a : x_1, b : x_2\}$ to
% the type $ \mathsf{record}\ \{ a : y_1, b : y_2\} $.
% But a substitution may later specialise $x_{1/2}$ and $y_{1/2}$ to incompatible
% layouts and types in various ways (including the overlapping of layouts), hence the need for remembering the constraint and checking it again
% when encountering a substitution.


% This rule is actually the only one that requires a constraint to be present in $C$.
% Note that such a constraint always relates a layout with a record type.

% The typing rule for specialising polymorphic functions now uses the
% specific notion of wellformedness for substitutions:
% \[
%   \inferrule
%   {
%     \FunTyEnv{f}{\forall \Delta'.C'
%       \Rightarrow \tau \rightarrow \tau'}  \\
%     {\Delta;C \vdash \varsigma,\delta :_s \Delta'; C'}
%   }
%   {\Delta ; C ;\Gamma \vdash f[\delta,\varsigma] : (\tau \rightarrow \tau')[\delta,\varsigma] }
%   {\textsc{Fun}}
% \]


%\TODO{The subtyping needs to be adapted to allow reordering the fields
%(or not, as we may assume they are already canonically ordered, since the
%formalisation supposedly accounts for the core syntax?)}

%\subsubsection{Future extensions}
%% Introducing this general set of constraints in the context could be avoided
%% if we restrict specialising layout variables to monomorphised types.
%
%% In this case, we could also reuse the already formalised simplified types
%% which do not have any variable case.
%
%% By exploiting the preservation of the typing rules by increasing the set of
%% constraints, it is straightforward to design a procedure that infers the necessary
%% set of constraints for a layout polymorphic function to be well-typed (if it can
%% be).
%
%Our type system currently does not check that the constraints
%in $C$ are compatible. For example, nothing forbids $C$ from containing both
%$(x\{o\},\ConsN{U8})$ and $(x\{o\}, \ConsN{U16})$, although in that case the layout variable 
%$x$ can never be fully instantiated, because of the wellformed rule for 
%substitution mentioned above.
%%
%Note that this does not allow unsafe behaviour, since C code generation only
%involves completely monomorphised layouts and types.
%%
%We leave for future work the design of a (possibly partial) validator that would
%deal with such situations and could raise a typing error as early as possible before
%instantiation.


\subsection{Compiling Records: Custom Getters and Setters}
\label{ssec:gsetters}

Currently, records are the only built-in types
in \cogent that can be boxed. We therefore use records as an example to
discuss how \dargent layouts influence the generated C code.

Without custom layouts, a \cogent record is directly compiled to a C struct
with as many fields as the original record: If
$T = \{a : A, b : B\}$, then
$\sem{ T }$ is a pointer type to
$\texttt{struct}\ \{ \sem{ A }\ a; \sem{ B }\ b; \}$,
 where $\sem{ \cdot }$ denotes the compilation of \cogent types to C types.


The \Dargent extension to the compiler relies on the observation that
we are free to choose
what a \cogent boxed record is compiled to
as long as we provide getters and setters for each field, since they
account for all the available \cogent operations on boxed records.
Assigning a custom layout
$\ell$ to a \cogent record $T$ results in a C type that we denote by
$\sem{ T }_\ell$, consisting of a C struct with a fixed-sized array
of words as a single field.
The implementation chooses the word size to be 32 bits, primarily because
most \cogent programs we develop are targeting 32-bit embedded systems, but this is not fundamental to the design and can be easily made 
configurable to any word size.
It is worth mentioning that this does not mean that a layout has to be a neat
multiple of 32 bits in size. It is absolutely valid to have a record
layout like $\DLRec{ \FieldTy{a}{1\byte}, \FieldTy{b}{3\bit} }$, 
11 bits in total. When this layout is embedded in another layout,
e.g.\ $\DLRec{ \FieldTy{x}{\DLRec{ \FieldTy{a}{1\byte}, \FieldTy{b}{3\bit} }},
               \FieldTy{y}{2\byte} }$, 
the remaining bits after the field $\FieldN{b}$ will be used by
the following field $\FieldN{y}$, without any implicit padding.


The getters and setters for each field are generated
according to the layout. If a top-level boxed record $T$ contains
a field $a : A$, the C prototypes are
\begin{align*}
    & \sem{ A } \texttt{ get\_a }(\sem{ T }_\ell \; t) ;
\\
    & \texttt{void set\_a }(\sem{ T }_\ell  \; t, \sem{ A } \; a) ;
\end{align*}
Note that $\sem{A}$, the return type of the getter (and similarly for
the second argument of the setter function) does not involve the
layout $\ell$. If $A$ is a boxed type, then $\sem{A}$ is a pointer to the
type $A$, whose layout is dictated by the layout information stored in
$A$'s sigil, independent of $\ell$. If $A$ is unboxed, the getter function is the point at which
we convert the low-level custom representation governed by the layout $\ell$
into a standard representation of $A$, so that the value of the field
$a$ can be inspected by the rest of the program. 
As an example, consider the $\FieldN{struct}$ field in \autoref{fig:cogent-type-ex}.
Roughly, the generated C getter has prototype
\cc|struct { U32 a; bool b; } get_struct (Example * t)|, where
$\AbsN{Example}$ is a C structure with a single field consisting of a fixed-sized array
spanning 16~bytes.
% Note that accessing the nested subfield
% \cg|a| or \cg{b} in \cogent involves first retrieving the whole top-level field
% \cg|struct|, which is inefficient and could be improved in future
% work, in particular since nested field getters and setters are already generated
% anyway, as auxiliary functions.


% \TODO{zilinc: describe the nested getters/setters more carefully
% to avoid the confusion that I had.}

% \subsubsection{Implementation}
% \review{Meanwhile, the paper spends a lot of space talking about things like layout
%   variables that seem neither very surprising nor very important. That space would
%   be used better to discuss, for example, the details of the generated getters and
%   setters.}

Getters and setters are generated recursively following the structure of the
involved field types. The process typically involves generating
auxiliary ``nested setters and getters'' that
directly manipulate the data array based on the value of the nested field (such as
$\FieldN{a}$ or $\FieldN{b}$ in the example of \autoref{fig:cogent-type-ex}), and similarly
for getters.
% For example, For nested unboxed records, such as \cg|struct| in
% Figure~\ref{fig:cogent-type-ex}, the compiler generates auxiliary ``nested
% setters and getters'' that
% directly manipulate the array based on the value of the nested field (such as
% $\FieldN{a}$ or $\FieldN{b}$), and similarly
% for getters.
 % Setting a field which is an unboxed record results in successive calls to these
 % nested setters.
For example, the getter and setter for the \cg|struct| field of 
\autoref{fig:cogent-type-ex} are roughly implemented as follows:
\begin{lstlisting}[language=C]
  // the data array
  typedef struct Example { U32[4] data } Example;

  struct field_struct { U32 a; bool b; };

  // setter for the struct field
  void set_struct(Example * d, struct field_struct v){
    // calls to nested setters
    set_struct_a(d, v.a);
    set_struct_b(d, v.b);
  } 

  // getter for the struct field
  struct field_struct get_struct(Example * d){
    // calls to nested getters
    return (struct field_struct){.a = get_struct_a(d), .b = get_struct_b(d)};
  } 
 \end{lstlisting}
% Getters follow a similar pattern.
% As for variant field, such as 
% Similarly, updating a variant field results in first updating the tag value in
% the data array, and then calling the generated nested setter for the payload.


% All getters and setters are directly accessing the top-level boxed type.
% In the example shown in Figure~\ref{fig:cogent-type-ex}, the getters for
% its child field $\FieldN{struct}$, grandchild fields $\FieldN{a}$ and $\FieldN{b}$ will all get from $\ConsN{Example}$, since the data structure is already
% flat.
\noindent As can be seen, a getter function (or similarly a setter) incurs a data format conversion:
It turns the low-level data format into a high-level typed value in C. Therefore
getting and passing around a large unboxed type can be expensive at run-time.
In practice though, if the program is carefully designed and implemented,
and the programmer only passes the minimal structures needed to external
functions, there is a good chance that an optimising C compiler is able to
eliminate unnecessary data conversions. \cogent allows programmers to
annotate functions with a \texttt{CINLINE} pragma, so that the \cc|inline|
modifier is generated in the C code, exposing more optimisation
opportunities. 
%In future, we plan to extend \cogent's uniqueness type system
%to allow for pointer access to unboxed types, improving the run-time 
%performance and reducing reliance on the C optimiser in these cases.


% As mentioned above, nested field getters and setters are generated 
% when nested (unboxed) records are involved, like \lstinline|struct| in Figure~\ref{fig:cogent-type-ex}.

\subsubsection{Invalid Bit Patterns}
Calling a getter on an external word array can lead to unexpected behaviours
when the
data format is invalid. This can happen when variant types are involved, if
the value in the tag part does not match any tag values declared in the
\dargent layout. Any undefined tag values will be treated as the last
constructor's tag. For example, if the layout of a variant
is $\DLVar{2\bit}{\ConsN{A}(0) : 1\byte, \ConsN{B}(1) : 2\byte,
\ConsN{C}(2) : 4\byte}$
and the tag value seen in the data format is 3, which does not match
any defined tag values but also fits in
the 2-bit space reserved for the tag, it will be deemed as the last
alternative, namely $\ConsN{C}$.

\dargent is not a low-level data parsing language, nor an interface language,
therefore the generated C getters do not check for validity, and
assume that the input is valid. Users of the language are responsible for
checking the validity of any incoming data.

\subsubsection{Required Properties of Getters and Setters}
As mentioned above, the compilation of a \cogent record type $T$ to C
can be arbitrary as long as getters and setters are provided. When it
comes to formally verifying the compiled C program, these getters and setters
are expected to compose well:
\begin{enumerate}\item 
  setting a field does not affect the result of getting another field;
\item getting a set field should return the set value (or at least an
  equivalent value).
\end{enumerate}
%
These properties are explained 
in~\autoref{sec:get-set-lemmas}. It is worth mentioning that, somewhat unintuitively,
extending the verification
framework to account for \Dargent does not require proving that
the generated getters and setters are accessing data at the locations
specified by the layout.

\subsection{Implementation}

As mentioned in~\autoref{ss:dargent-core}, the \dargent core calculus initially represents a bit range
as a pair of integers denoted by $\bitrange{o}{s}$, where $o$ indicates the offset (in bits)
to the beginning of the top-level datatype in which it is contained, and $s$ indicates how many
bits it occupies. This representation is concise and easy to work with when typechecking the core language.
Such a representation, however, does not necessarily lend itself to an easy code
generation algorithm. Therefore we have another
step to convert each bit range into a list of \emph{aligned bit ranges} tailored for C code generation. Each aligned bit range is essentially a word.

An aligned bit range $\abr{w}{o}{s}$ is a triple of integers, where $w$ is the
word-offset to the top-level datatype, $o$ is the bit-offset inside this word,
and $s$ is the number of bits occupied.
Each aligned bit range contains the information of how the bits in each word is used.
The generated C getter/setter consists of a series of statements, each operating on one
word via bitwise operations.
Each such C statement is generated according to one aligned bit range. This one-to-one mapping simplifies
C code generation. Because of the deployment of aligned bit ranges,
it is easy for us to choose an appropriate word size (namely the size of the
aligned bit ranges) specifically tailored for the application in question.
This flexibility is important for low-level programming (see \autoref{ssec:power-control-system} for a concrete example) and for high performance.
% The intuition is that machine instructions typically work on the granularity of machine words.
% In the generated C code, we want to mimic that.
% have to partition the layout encoded in a bit range into words, each represented as
% an aligned bit range, so that they closely correspond to the C statements we are going to generate.
% ~\footnote{This is with the exception of array compilation, which we will cover shortly.

% For this reason, we break down a bit range into smaller units, thus each one can be operated on by
% a single statement in C.
% The granularity, we call \emph{aligned size}, is thus set to be a machine word.
% (but in general can be different, as we shall see when coming to arrays).
% One key meta-property that holds on aligned bit ranges is as follows:
% $$\forall \abr{w}{o}{s}.\, w \ge 0 \wedge o \ge 0 \wedge s \ge 1 \wedge w + o \le A$$
% where $A$ is the align size.

% The compilation of arrays is handled slightly differently. In the current language
% design, we restrict the layout of each element to be byte-aligned. It means that
% each element will have to be at least one byte long. We thus set the granularity to
% a byte, rather than a machine word, when we encounter arrays.


